\documentclass[11pt]{article}
\usepackage{amsmath,amssymb}
\usepackage{longtable,array}
\usepackage[margin=1in]{geometry}
\usepackage{hyperref}

\title{Ideas Register --- Consolidated, v1 to v20}
\author{One file for every parked idea, with its current status}
\date{}

\begin{document}
\maketitle

\noindent This file replaces \texttt{ideas\_v1} through
\texttt{ideas\_v18}. Versions 1--11 were cumulative (each carried the
previous forward) and 12--18 were separate additions, so the same
idea often appeared in several files at different stages. Here each
idea appears once, with where it came from, where it is destined, and
where it stands now.

\textbf{How status was determined.} Not from memory of the
discussions. Each status was checked by searching the current paper,
\texttt{Pattern\_Arithmetic\_v15\_22.tex}, and --- for items destined
elsewhere --- the \texttt{Assembly\_v1.tex} and
\texttt{Inference\_v1.tex} copies on hand (which may be older than
your working copies; if so, the ``not filed'' marks below should be
rechecked against yours). Entries from v19 onward are appended here; v19 and v20 are the first two.

\paragraph{Status vocabulary.}
\textbf{CLOSED} --- resolved and merged into the paper.
\textbf{CLOSED (redirected)} --- determined not to be the paper's
question; destination named.
\textbf{RESOLVED, NOT MERGED} --- fully worked out but waiting on a
prerequisite.
\textbf{DECIDED} --- a ruling recorded; nothing to merge.
\textbf{PARTLY MERGED} --- some of it is in the paper.
\textbf{OPEN} --- still unresolved or unfiled.
\textbf{PARKED} --- verified but with no current home.

\section{The gaps of ideas v1, as they now stand}

Honest summary first, since this is what prompted the consolidation:

\begin{center}
\begin{tabular}{p{4.2cm}p{10.4cm}}
Gap & Status \\
\hline
\textbf{Gap 1} --- Mix across differing directions &
\textbf{CLOSED} in \texttt{v15\_21}: refused, on the $\mathtt{chladni}$
evidence, with Bundle and Collapse remaining available. One small
residual (R26 below). \\[4pt]
\textbf{Gap 2} --- Location as a sort &
\textbf{RESOLVED, NOT CLOSED.} Fully worked out in v17 (meta vector;
both encodings verified). It cannot honestly be marked closed in the
paper yet: it depends on the meta-vector taxonomy being imported into
Paper~1 first, and that taxonomy is not in the paper at all (zero
occurrences). The paper's own ``Open: location as a sort'' paragraph
is therefore still correct as written. \\[4pt]
\textbf{Gap 3} --- Reducing a path by merging close actors &
\textbf{CLOSED for Paper~1} in \texttt{v15\_22} (narrowed to the
algebraic guarantee); the technique question is \textbf{redirected to
Inference} and is not yet filed there. \\
\end{tabular}
\end{center}

The two remaining items from v1 are both Inference-destined and
neither is filed: the candlestick summary (verified, mechanism
works) and the ``Inference Mandelbrot'' naming flag.

\textbf{The single most useful next step} is importing the
meta-vector taxonomy into Paper~1 (R15). It is the prerequisite for
closing Gap~2, and it gives a home to location, the two time
registers, and $V_{\text{derived}}$ at once.

\section{Register}

\small
\setlength{\tabcolsep}{4pt}
\newcolumntype{R}[1]{>{\raggedright\arraybackslash}p{#1}}
\begin{longtable}{@{}l R{4.9cm} R{1.5cm} R{2.3cm} R{4.4cm}@{}}
ID & Item & Origin & Destination & Status \\
\hline
\endhead

\multicolumn{5}{@{}l}{\textit{Closed}} \\*
R1 & Mix across differing directions (Gap 1) & v1, v3--v4, v7 & Paper 1 & \textbf{CLOSED} --- refused; v15\_21 \\
R2 & Terminology realignment: Bundle/Remove, Collapse/UnCollapse, Mix/UnMix & v8--v10 & Paper 1 & \textbf{CLOSED} --- merged; v15\_20 \\
R3 & Symbols $\uplus,\setminus,\oplus,\ominus,\boxplus,\boxminus$ & v9--v10 & Paper 1 & \textbf{CLOSED} --- merged ($\circleddash$ superseded) \\
R4 & Collapse non-idempotent; $N$ = list length; $A^{\langle N\rangle}$; UnCollapse algebra & v10 & Paper 1 & \textbf{CLOSED} --- merged \\
R5 & Distinct count $D$ against total $N$ & v13 & Paper 1 & \textbf{CLOSED} --- merged; v15\_20 \\
R6 & Reducing a path (Gap 3): algebraic guarantee only & v1, v6, v18 & Paper 1 & \textbf{CLOSED (redirected)} --- v15\_22 \\
R7 & $\tau_{\text{weather}}$ conflation & v4 & --- & \textbf{CLOSED} --- correction recorded \\
R8 & Same-sort region merge vs cross-sort & v6 & Assembly; Paper 1 & \textbf{CLOSED} --- same-sort solved in Assembly; cross-sort is R1 \\
R9 & Internal link, terminology, cohesion audit & v12 & Paper 1 & \textbf{CLOSED} --- fixes in v15\_23, carried forward \\
R10 & Stale status lines and symbols in v1--v10 & v11 & ideas & \textbf{CLOSED} --- superseded by this file \\
R11 & Inertia in Assembly via a mass-like weighting & v13 & Paper 2 & \textbf{CLOSED} --- tested; fails \\[4pt]

\multicolumn{5}{@{}l}{\textit{Decided}} \\*
R12 & Physics names for atom metrics & v13 & Paper 1 & \textbf{DECIDED} --- keep $\mathcal{E}$; $N$/$D$ earned; ``density'' stays retired; $S/N$ parked \\
R13 & ``uncollapse'' (v5) collides with UnCollapse & v11 & ideas & \textbf{DECIDED} --- call it Mix-provenance link \\[4pt]

\multicolumn{5}{@{}l}{\textit{Resolved but not merged, and its prerequisite}} \\*
R14 & Gap 2: location and time as meta vectors; registers; age derived & v1, v17 & Paper 1 & \textbf{RESOLVED, NOT MERGED} --- blocked on R15 \\
R15 & Meta-vector taxonomy; $V_{\text{derived}}$ & v2 & Paper 1 & \textbf{OPEN} --- not in paper; prerequisite for R14 \\[4pt]

\multicolumn{5}{@{}l}{\textit{Open --- meta vectors}} \\*
R16 & Link vectors & v2 & Paper 2 & \textbf{OPEN} --- not filed \\
R17 & Meta-vector grouping (promote a shared one) & v2 & Inference & \textbf{OPEN} --- not filed \\
R18 & Identity handle / Mix-provenance link & v5, v11 & Papers 1, 2, Inf. & \textbf{OPEN} --- paper has ``proposed, not adopted'' only \\
R19 & Should Mix/Collapse skip expired contributions? & v17 & Paper 1; Inf. & \textbf{OPEN} --- deliberately unanswered \\[4pt]

\multicolumn{5}{@{}l}{\textit{Open --- Collapse}} \\*
R20 & Collapse across differing sorts (per-vector $\tau_i$) & v10, v11 & Paper 1 & \textbf{OPEN} --- stated in paper; broadened by R36 \\
R21 & Collapse as $T$'s actor & v11 & Paper 1 & \textbf{OPEN} --- ``not claimed'' in Axiom 7; see R40 \\
R22 & Assembly and Collapse: relation & v11 & Papers 1, 2 & \textbf{OPEN} --- Axiom 6 silent \\[4pt]

\multicolumn{5}{@{}l}{\textit{Open --- metrics, scope, presentation}} \\*
R23 & Shell-indexed energy & v14 & Paper 1 & \textbf{OPEN} --- undefined \\
R24 & Antipodal ($\hat n$ against $-\hat n$) cancellation & v14 & Paper 1 & \textbf{OPEN} --- special case of closed R1 \\
R25 & Separable-metric pointwise test & v15 & Paper 1 & \textbf{PARKED} --- verified; value depends on R27 \\
R26 & Route-1 scoping clause (refusal is about $\boxplus$ only) & v7 & Paper 1 & \textbf{PARTLY MERGED} --- one sentence not carried \\
R27 & Axiom 7 metric catalogue and costumes: cut or move? & v16 & Paper 1 & \textbf{OPEN} --- undecided \\
R28 & Two-paper split & v16 & Papers 1, 2 & \textbf{OPEN} --- undecided \\
R29 & Review's presentational items & v16 & Papers 1, 2 & \textbf{OPEN} --- not acted on \\[4pt]

\multicolumn{5}{@{}l}{\textit{Open --- Inference-bound}} \\*
R30 & Candlestick summary & v1 & Inference & \textbf{OPEN} --- verified; not filed \\
R31 & ``Inference Mandelbrot'' naming flag & v1 & Inference & \textbf{OPEN} --- flag only \\
R32 & Technique selection for reduction (average, centroids, shape-preserving) & v18 & Inference & \textbf{OPEN} --- specified; not filed \\
R33 & Douglas--Peucker adapted to the sphere & v18 & Inference & \textbf{OPEN} --- the one unbuilt piece \\
R34 & Sphere-packing bound for the ``stated distance'' & v6 & Inference & \textbf{OPEN} --- not filed anywhere \\
R35 & Coherence check for the candlestick & v1 & Inference & \textbf{OPEN} --- never considered \\[4pt]

\multicolumn{5}{@{}l}{\textit{Open --- atom structure (v19)}} \\*
R36 & Pure and mixed atoms: cores; per-vector $\tau_i$ as the general form; Mona Lisa as test case; direction-based derived quantities per sort & v19 & Paper 1 & \textbf{OPEN} --- broadens R20 \\
R37 & Collapse of multi-vector contributions (not fixed at one vector each) & v19 & Paper 1 & \textbf{RESOLVED, NOT MERGED} --- paper still says one vector per contribution \\
R38 & $N$ as number of cores; $Q$ as number of vectors (resolves the collision with $N$ as vector count in Axiom 7) & v19 & Paper 1 & \textbf{DECIDED} --- $N$ cores, $Q$ vectors; paper edits wait for the merge \\
R39 & Collapsed atom as a counted bowl; commuting test & v19 & Paper 1 & \textbf{OPEN} --- test not yet run \\
R40 & $T$ asymmetry: a bowl is not an actor; a Collapsed atom needs a cast & v19 & Paper 1 & \textbf{OPEN} --- relates to R21 \\[4pt]

\multicolumn{5}{@{}l}{\textit{Open --- energy and effective counts (v20)}} \\*
R41 & Energy from strength, vector count and direction: $\mathcal E=S^2/Q+Q\,\mathrm{Var}(r)$; $S^2/Q\le\mathcal E\le S$; $|V|^2=(1-\alpha)\mathcal E+\alpha S^2$ & v20 & Paper 1 & \textbf{PARKED} --- identities verified, not in paper; $\mathcal E$ keeps ``energy,'' $\mathcal V$ keeps ``volume'' \\
R42 & Effective vector count $Q_{\mathrm{eff}}=S^2/\mathcal E$; overlap with $\sigma_{\text{atom}}$ & v20 & Paper 1 & \textbf{OPEN} --- candidate; keep it, $\sigma_{\text{atom}}$, both, or neither \\
R43 & Strength-only metrics per sort and for the whole; effective number of sorts $1/\sum w_s^2$; pooling ratio & v20 & Paper 1 & \textbf{OPEN} --- needs one stated convention \\
\end{longtable}
\normalsize

\section{The entries, by version}

\subsection*{v1 --- the three gaps and two Inference items}
\textbf{Gap 1} (Mix across differing directions): Axiom~5 defines Mix
only at a shared $\hat n$; a provisional vector-sum recipe had been
tried and withdrawn. \emph{Status: R1, closed.}
\textbf{Gap 2} (location as a sort): a third candidate was tested ---
disjoint dimensions for content, time, and location. Verified: two
shifts through disjoint dimensions commute exactly (to full
floating-point precision), and content stays recoverable by
projection; shifts forced onto one shared dimension do not commute.
Left open: the encoding itself. \emph{Status: R14, resolved in v17.}
\textbf{Gap 3} (reducing a path): needed a distance, a merge rule,
and a drift bound. Verified that swapping the clip for the Hill form
$r'=|z|/(K+|z|)$ changes the strength but not $\chi'=\arg(z)$.
\emph{Status: R6, closed (redirected).}
\textbf{Candlestick summary}: open/close/high/low distinguishes two
sequences of identical content in opposite order that naive Mix cannot
--- Mix gave the same $47.5^\circ$ for both; the candle's open and
close swapped ($20^\circ\!\to\!160^\circ$ against
$160^\circ\!\to\!20^\circ$) while high and low stayed equal.
\emph{Status: R30, 0 occurrences in \texttt{Inference\_v1}.}
\textbf{``Inference Mandelbrot''}: an ordered-thing promotion is a
different problem from Assembly's Hill-bijection promotion of
unordered camps, and needs its own name. \emph{Status: R31.}
The ``considered and found to need nothing filed'' list (reference
atoms, precursor-witness relabeling, the $\pi=(x;w_1,\dots,w_n)$
trajectory, $\chi$ as internal phase-color) was rechecked in v11
against the current paper: all present, nothing to file. \emph{Closed.}

\subsection*{v2 --- meta vectors, split across three papers}
Proposed: besides information vectors (capped, $r_i\le1$), an atom
may carry meta vectors describing it rather than being part of its
content: $V_{\text{derived}}$, link vectors, a timestamp vector.
Meta vectors are not capped. Split by destination: the taxonomy and
$V_{\text{derived}}$ to Paper~1; link vectors to Assembly; the
grouping mechanism (query nearby meta vectors, promote one shared
value) to Inference. Verified caveat: grouping $V_{\text{derived}}$
inherits the Hill form's limit --- both fix magnitude only, never
direction. \emph{Status: R15--R17. The paper has zero occurrences of
``meta vector''; \texttt{Assembly\_v1} has zero of ``link vector.''}

\subsection*{v3 --- Gap 1, characterized}
Argued reading: Mix takes several reports already agreed to answer
one shared question and returns a verdict about it; the shared
$\hat n$ is what makes that true. Notably, v3 already stated that in
Mix's domain \emph{``$\hat n$ itself was never plural; only $\chi$
ever varied''} --- the very distinction that had to be rediscovered
during the Gap~3 discussion (v18) before it stopped causing
confusion. Four ``we mix all the time'' counterexamples were checked
and dissolved (sugar in tea is one dial; red plus green is already
same-$\hat n$; paint compositing is a bowl and flattening is
destructive and irreversible --- matching UnMix's impossibility). The
one unresolved example, a whole painting as one vector, is Gap~1
itself. The open question split in two: \emph{embedding validity}
(does closeness on a shared sphere track semantic agreement?) and the
\emph{merge rule}. \emph{Status: superseded by R1's closure --- Mix
refuses, so neither sub-question needs an answer for Mix.}

\subsection*{v4 --- the $\tau_{\text{weather}}$ correction; Gap 1 ``ready to close''}
$\tau_{\text{weather}}$ had named two different problems: combining
different weather concepts across directions (Gap~1), and reducing
closely aligned same-sort atoms to one representative, which is
Assembly's already-verified Hill-bijection mechanism (16 camps to
13 systems). The $\mathtt{chladni}$ section was found to already
contain the evidence for Gap~1 (``cross-mode coexistence is
therefore $\uplus$''). \emph{Status: R7 closed; the $\mathtt{chladni}$
evidence is what the Gap~1 closure rests on.}

\subsection*{v5 --- identity handle}
An opaque $id$ meta field, assigned before Bundle or Mix and kept
after, so a result can point back to the bowl that produced it
(bowl $\{\text{red},\text{green}\}$ is $\#17$; Mix gives yellow
$\#18$ with a link $\#18\to\#17$). Refuse $\tau_{id}$ as a sort; keep
$id$ as a field. Navigation by a kept record, not inversion of Mix.
Renamed in v11 ``Mix-provenance link'' because the original word
``uncollapse'' collides with the formal UnCollapse operator.
\emph{Status: R13 decided; R18 open --- the paper mentions $id_i$
once, as ``proposed, not adopted,'' in UnCollapse's well-definedness
paragraph.}

\subsection*{v6 --- sphere-packing bound; merging a whole region}
For Gap~3's ``stated distance'': the number of mutually
distinguishable points at separation $\theta_{\min}$ is roughly
$2/(1-\cos(\theta_{\min}/2))$ --- about $52{,}525$ at $1^\circ$,
$2{,}101$ at $5^\circ$, $526$ at $10^\circ$, $59$ at $30^\circ$, $15$
at $60^\circ$, $7$ at $90^\circ$, $4$ at $120^\circ$. Merging a
whole region to one representative splits: same-sort is solved
(Hill bijection) and cross-sort is Gap~1. \emph{Status: R8 closed;
the bound itself is R34, filed nowhere.}

\subsection*{v7 --- scoping of Route 1}
``Refuse outright'' should be worded as a claim about $\boxplus$,
not a permanent bar on any future, separately admitted mechanism ---
the same ``undecided, not admitted'' logic Axiom~0 already uses.
\emph{Status: R26, partly merged.} Checked against the Gap~1
closure: the closure is scoped to Mix and keeps $\uplus$ and
$\oplus$ available, which honours the point. What was not carried is
the one explicit sentence that a separately admitted mechanism is not
barred, and the closure calls the other two routes ``withdrawn,''
slightly stronger than v7 intended. An optional one-sentence fix.

\subsection*{v8, v9, v10 --- terminology and verb inventory}
v8 proposed the Collapse/UnCollapse/Bundle/Mix/UnMix realignment and
flagged a hard symbol collision ($\ominus$ was Remove). v9 chose
$\uplus$ for Bundle and $\circleddash$ for UnCollapse, and scoped
ids to the one case that needs them: duplicate-valued contributions
in UnCollapse. v10 revised the symbols --- $\ominus$ to UnCollapse,
Remove to $\setminus$ --- settled Bundle's idempotence
($A\uplus A=A$, already proven), and showed Collapse is
\emph{not} idempotent by three convergent arguments (information
loss; an algebraic contradiction with $A\ \text{UnCollapse}\ A=0$;
the same evidence-count blindness the Hill form was built to fix).
$N$ is simply the length of the vector list; $A^{\langle N\rangle}$
abbreviates $N$ identical copies; UnCollapse sends $N\to N-1$.
\emph{Status: R2--R4 closed, all merged. MixMerge was later reverted
to Mix at your request; $\circleddash$ appears zero times in the
paper. The per-vector $\tau_i$ question that v10 left open is R20.}

\subsection*{v11 --- review of v10 against the paper}
Status lines of v8--v10 corrected (the realignment was merged); stale
$\oplus$/$\ominus$ readings in earlier text fixed, including one
sentence that had become \emph{false} (``Gather still satisfies
$A\oplus A=A$'' reads wrongly once $\oplus$ means Collapse); the
``uncollapse'' collision resolved by renaming. Three gaps opened by
giving Collapse axiom-level standing: R20 (across sorts), R21 (as an
actor for $T$), R22 (no stated relation to Assembly --- Axiom~6
mentions Collapse zero times, so this one is silence rather than a
stated omission). One new route recorded: Collapse losslessly first,
then define a reading --- the seed of what v18 developed.
\emph{Status: R10, R13 closed or decided; R20--R22 open.}

\subsection*{v12 --- internal audit}
96 labels and 198 reference call-sites checked, each against what
its sentence claims; zero mismatches. Two stragglers fixed (an
Abstract ``a Collapse'' noun phrase; a Forest-costume ``gathers'').
Four cohesion fixes where the paper's own summaries omitted Collapse,
and the UnMix proof section now names UnMix. No new conceptual gap.
\emph{Status: R9 closed; fixes are in the current paper.}

\subsection*{v13 --- naming the atom metrics}
Seven metrics, quoted precisely. Earlier drafts had called
$\mathcal{V}_{\text{atom}}$ ``mass,'' $\gamma_{\text{atom}}$
``density,'' and $S_{\text{atom}}$ ``atom energy,'' and retired all
three; only $\mathcal{E}_{\text{atom}}=\sum r_i^2$ kept ``energy,''
because it is tied by an exact identity to
$|V_{\text{derived}}|^2=\mathcal{E}_{\text{atom}}+2\sum_{i<j}r_ir_j(\hat n_i\!\cdot\!\hat n_j)$
(variance-of-a-sum structure, verified to six decimals). The rule:
a physics name is kept only if the quantity participates in something
provable. $D=|A_1\uplus\cdots\uplus A_n|$ passes (forced by $\uplus$'s
idempotence), so $N$ and $D$ keep the mass-number/atomic-number
analogy; ``density'' stays retired; $S_{\text{atom}}/N$ is parked
rather than named. Inertia in Assembly is closed: a mass-like
weighting was already tested in the motion equation --- a ninefold
strength gap gives a $729$-fold weight and a non-monotonic
trajectory.
\emph{Added in conversation after v15, not previously filed:} asked
whether either of the two metrics survives, the answer was that
energy already does ($\mathcal{E}_{\text{atom}}$, earned) while no
candidate for density does --- $S_{\text{atom}}$ already has a name,
$N$ is a count rather than a ratio, and the one ratio with the right
shape, $\gamma_{\text{atom}}$, is the one already rejected.
\emph{Status: R5 closed; R11 closed; R12 decided.}

\subsection*{v14 --- three rough sketches}
\emph{Density} restated: nothing new, closed by reference to v13.
\emph{Energy as shell distribution}: partly already true of
$\mathcal{E}_{\text{atom}}$ --- for fixed total strength $S=2.0$,
four vectors give $\mathcal{E}=1.0$ spread evenly and $1.81$
concentrated toward the extremes (and $\sigma_{\text{atom}}$ goes
$0\to0.333$) --- but $\sigma_{\text{atom}}$ is blind to \emph{which
side} the concentration is on, and no shell-indexed quantity exists.
\emph{Matter/antimatter}: antipodal cancellation ($\hat n$ against
$-\hat n$), a second kind of opposition beside the existing
same-direction opposite-phase cancellation. With Gap~1 now closed by
refusal, this would be a deliberate special-case exception, parallel
to the existing one; what the rule would produce, and whether it
needs equal strength, is undecided.
\emph{Status: R23, R24 open.}

\subsection*{v15 --- when two metrics constrain each other}
For separable metrics $M_1=\sum f(r_i)$ and $M_2=\sum g(r_i)$: if
$f\le g$ everywhere on $[0,1]$ then $M_1\le M_2$ for every atom;
if the curves cross there is no bound. Verified against both cases
(it forces $\mathcal{E}\le S$ because $r^2\le r$; it forces nothing
for $N$ against a shell-weighted energy). Where a bound holds, the
gap is itself a metric ($r-r^2$, largest at $r=0.5$). Possible uses
(a non-redundant reporting set, a two-axis placement, a debugging
check, a post-hoc label for Assembly's camps) are speculative.
\emph{Status: R25, parked; it matters mainly if R27 keeps the
metrics.}

\subsection*{v16 --- external review and the two-paper split}
Every checkable claim in the review was verified: the worked
numbers (e.g.\ $r'\approx0.721$ at $73.9^\circ$; freeze at $628$ and
$1261$ steps), that all 22 occurrences of the six Axiom~7 metrics sit
inside Axiom~7 and none in Bundle/Mix/$T$/offset, that ``white'' is
used in two senses (the paper already says so in Axiom~3), and that
the ``checked directly'' precision claims have no supplement.
The proposal: Paper~1 is the typed algebra, Paper~2 is assembly and
the enclosure, the Axiom~7 catalogue and the costumes go in neither;
Paper~1 states Assembly as a signature only and ends by naming the
seat it refuses to supply; cut by question, not by length.
\emph{Not itemized when v16 was filed, but in the review and still
pending (R29):} state the clip, the Mix refusal, and the assembly
energy as stipulated choices in one place rather than narrated as
discoveries; print the enclosure gradient; provide a supplement for
the ``checked directly'' claims; and stop calling the enclosure a
``law'' in the same voice as $\mathrm{SO}(3)$ while it still needs
$\delta$, $\eta$, the census cut, and $\varepsilon_{\text{allow}}$.
\emph{Status: R27, R28, R29 open and undecided.}

\subsection*{v17 --- Gap 2 resolved}
Meta vector, not sort, tested against the cap: $r\le1$ encodes
strength of content, while location and time are coordinates. This is
the line the paper already draws for Assembly's seat. Location is
one vector via standard $(\theta,\phi)$ from latitude and longitude
(London $\to(0.6225,-0.0011,0.7826)$, Tokyo
$\to(-0.6194,0.5252,0.5835)$, Sydney $\to(-0.7273,0.3999,-0.5577)$).
Time on $(\theta,\phi)$ fails: at hour $0$ every minute lands on
$(0,0,1)$, and $11{:}59$ and $0{:}00$ land about $2$ apart
(antipodal). The fix: minute on $\phi$, hour on $\chi$, $\theta$
fixed at the equator ($6^\circ$ and under $1^\circ$ apart for the
same pair). Timestamp is a register of length $N$, kept in step with
UnCollapse; age is derived ($\text{now}-\text{timestamp}$), never
stored; expiry is a second register because it cannot be derived;
timestamps help but do not guarantee disambiguation the way an $id$
would. \textbf{Blocking:} cannot be merged until R15 is.
\emph{Status: R14, R19.}

\subsection*{v18 --- Gap 3 closed for Paper 1}
Corrected an earlier confusion ($\chi$ is phase at a shared direction,
not $\hat n$; Gap~1's refusal did not apply). The real issue is
order-blindness: merging three actors in either order gives the same
result. A path's commitment to order comes from outside the algebra
--- Bundle discards order, and choosing $T$ over a bowl is the
declaration that order matters. So averaging 10{,}000 temperature
readings is bowl-plus-Mix and loses nothing. Three techniques answer
three different questions: averaging cannot tell a spike from a flat
line; centroids give identical results for a sequence and its reverse
(clustering is blind to position); shape-preserving reduction keeps
real points at their positions --- Douglas--Peucker, whose tolerance
is the drift bound Gap~3 asked for. Paper~1 was narrowed to the
non-consumption guarantee. \emph{Status: R6 closed (redirected);
R32, R33 open for Inference.}


\subsection*{v19 --- pure and mixed atoms, cores, and Collapse against a bowl}
Prompted by a hypothetical: could the Mona Lisa be held in one atom?
``One atom'' can mean three things. \emph{A name}: a single seated
point, like ``storm,'' which the paper calls a doorbell --- it carries
none of the content. \emph{A bowl}: already legal for mixed sorts,
since Bundle is sort-blind. \emph{A Collapsed atom} holding colors,
words, and strokes, which needs mixed sorts inside one atom. This is
the example v3 left unresolved (``a whole painting as a single
vector'' was Gap~1 restated). With Gap~1 closed, a single
\emph{vector} is refused, while one atom of many vectors stays
available. What stands in the way of the third meaning: mixed sorts
(R20, widened to R36); where each contribution sits on the canvas ---
a location per contribution, the same kind of register as v17's but
not yet specified per contribution; derived quantities across sorts;
and relations, since the paper ``does not mint the relations'' and,
by its own design rule, which units were co-present and within what
region are facts of a level, not coordinates of an atom.

\pagebreak[4]
\textbf{Rulings given in the discussion, none yet in the paper.}\par\nopagebreak
\begin{itemize}
\item \emph{Collapse juxtaposes.} Positions do not change, nothing
merges, no new type is minted. Relations within a sort are as before;
across sorts they exist only under a stated rule (Axiom~0's
``undecided, not admitted,'' applied to Collapse). So derived
quantities that use direction are computed per sort, and cross-sort ones
are undefined until a rule exists. (Refined in v20: strength-only
quantities may also be given for the whole atom.)
\item \emph{Cores.} A core is a unit inside an atom, usually one
$\tau$ (a pure type) but free to hold several --- words and numbers
together --- as typical, not as a rule. The general form is
per-vector $\tau_i$, with ``core'' a grouping. ``Core'' appears in
the paper only in the ``Core Contributions'' heading and in ``no
core, no halo'' in the enclosure, so it needs one defined sense.
\item \emph{Collapse is not fixed at one vector per contribution.}
This settles the question the paper was silent on: whether collapsing
already-multi-vector atoms flattens or nests. The Collapse text still
says ``one vector per contribution.''
\item \emph{$N$ is the number of cores.} This keeps UnCollapse's
$N\to N-1$, keeps $D$ a bowl size (distinct cores), and keeps the v17
registers at length $N$, so those registers are per core, not per
vector. The cost is a word collision: $N$ is currently the vector
count in Axiom~7, in $\gamma_{\text{atom}}=\mathcal V_{\text{atom}}/N$
and in $\sigma_{\text{atom}}$'s $N^2/(N-1)$. \emph{Decided:} $N$ counts cores and $Q$ counts vectors. All seven
Axiom~7 metrics run over vectors, so they take $Q$; the $D/N$
paragraph and the UnCollapse text keep $N$ and change ``vectors'' to
``cores.'' To settle at the merge: the sentence that defines $N$ as the
atom's count of informational vectors, and where the label ``info
number'' goes.
\item \emph{A Collapsed atom matches a counted bowl.} They are
identical exactly when nothing repeats ($D=N$), and $N-D$ is the gap.
The paper's bowl is a set, and its own counting equation shows the
difference: $A\uplus A\uplus A\uplus A\setminus A=\emptyset$, while
$A\oplus A\oplus A\oplus A\ominus A=A^{\langle3\rangle}$ (by the
UnCollapse rule). Operations should work on both; recorded as a
commuting test --- a pack/unpack map between the two, with every
operation required to agree across it. Sum-style quantities
($S$, $\mathcal E$, $V_{\text{derived}}$) agree only on a counted
bowl. $T$ is asymmetric today: it ``does not act on the set $R$ as a
set,'' and a Collapsed atom can act only once a cast is named.
\item \emph{Two distinct ways of arranging and viewing vectors.}
They differ in multiplicity (set against counted), in handle (a
Collapsed atom is one unit that can sit in a bowl or at a lexicon
pin; a bowl is a collection), and in sequence (neither stores order;
the path carries it, and a bowl is a good surface for displaying one).
\end{itemize}
\emph{Status: R36--R40; R20 and R21 cross-referenced.}


\subsection*{v20 --- energy from strength, counts, and direction; effective counts}
Started from a recollection that energy is $\sum r^3$. The paper's
energy is $\mathcal E=\sum r^2$; $\sum r^3$ is the atom's volume,
formerly ``mass.'' Nothing needs renaming, and $\sum r^3$ keeps a
geometric reading (strength as a radius, as in Assembly's softening
length).

\textbf{Energy from strength, vector count $Q$, and direction already
exists in $\mathcal E$.} Three exact relations, verified on 20{,}000
random atoms (errors near $10^{-14}$, no violations), with
$S=\sum r_i$:
\begin{itemize}
\item $\mathcal E=S^2/Q+Q\,\mathrm{Var}(r)$: an evenly spread floor,
set by $S$ and $Q$ alone, plus concentration.
\item $S^2/Q\le\mathcal E\le S$.
\item $|V_{\text{derived}}|^2=(1-\alpha)\mathcal E+\alpha S^2$: the
resultant's energy, blending the incoherent baseline with the fully
aligned value, weighted by the paper's coherence $\alpha$.
\end{itemize}
Which works better: the floor $S^2/Q$ cannot tell an even atom from a
concentrated one (both give $1.000$ at $S=2$, $Q=4$, where $\mathcal E$
gives $1.000$ against $1.810$). $\mathcal V=\sum r^3$ separates them even
more sharply ($0.500$ against $1.715$) but has no identity tying it to
anything beyond the radius reading. So $\mathcal E$ is the energy and
$\mathcal V$ stays volume. Cores $N$ enter through the same
variance-of-a-sum algebra applied per core (plain algebra, not run
separately) and through the ratios $Q/N$ and $S/N$. The mean strength
$S/Q$, the parked candidate from v13 (then written $S/N$), is the centre
of the variance decomposition and keeps a plain name.

\textbf{Effective vector count} $Q_{\mathrm{eff}}=S^2/\mathcal E$
(the participation ratio in physics, the effective sample size in
statistics).
\begin{itemize}
\item $Q$ counts vectors and $Q_{\mathrm{eff}}$ counts those that matter:
100 vectors at $0.05$ give $100$; one vector at $0.9$ plus 99 at
$0.001$ give $1.23$.
\item The largest vector carries at least $1/Q_{\mathrm{eff}}$ of the
strength (50{,}000 random atoms, no violations).
\item Scale-free: halving every $r$ leaves it unchanged ($2.2099$ both).
\item With direction: $|V|^2/S^2=\alpha+(1-\alpha)/Q_{\mathrm{eff}}$
(error near $7\times10^{-16}$), so two dimensionless numbers fix the
resultant's size.
\item Limits: it re-expresses $S$ and $\mathcal E$, adding no
information, and ignores direction.
\item Overlap: $\sigma_{\text{atom}}=(Q/Q_{\mathrm{eff},\mathcal V}-1)/(Q-1)$,
where $Q_{\mathrm{eff},\mathcal V}$ is the effective count of the
\emph{cubed} shares (error near $10^{-15}$). So $\sigma_{\text{atom}}$
is already this family, and $Q_{\mathrm{eff}}$ is its plain-strength
sibling, the one tied by identity to $\mathcal E$ and $|V|^2$. Whether
to keep both, one, or neither is undecided.
\end{itemize}

\textbf{Per sort and for the whole.} The per-sort rule belongs to
quantities that use direction ($V_{\text{derived}}$, $\alpha$,
$|V|^2$), since sums and dot products of directions across sorts mean
nothing. Strength-only quantities ($S$, $\mathcal E$, $\mathcal V$,
$Q_{\mathrm{eff}}$, the $\sigma$'s) can be given per sort and for the
whole, on one stated convention: each sort's $r$ is normalized to its
own declared saturation (the paper's $r=a/a_{\max}$), and whole-atom
figures treat fraction of own saturation as a common scale. With
$w_s$ the share of total strength held by sort $s$:
\begin{itemize}
\item Pooling ratio $\sum_sQ_{\mathrm{eff},s}/Q_{\mathrm{eff,whole}}\ge1$
(Cauchy--Schwarz; no violations in 30{,}000 atoms), with equality
exactly when every sort has the same strength-weighted mean strength
$\mathcal E_s/S_s$. It measures imbalance between sorts: two vectors at
$0.9$ against ten at $0.1$ give $2.63$.
\item Exact decomposition:
$1/Q_{\mathrm{eff,whole}}=\sum_sw_s^2/Q_{\mathrm{eff},s}$ (error near
$2\times10^{-16}$). In the symmetric case it multiplies: three sorts of
four vectors at $0.5$ give $12=3\times4$.
\item Effective number of sorts $1/\sum_sw_s^2$ answers ``is this atom
really mixed?'' more readably than the ratio: for 50 color vectors at
$0.5$, 3 word vectors at $0.5$, and 10 stroke vectors at $0.3$,
$Q_{\mathrm{eff}}$ is $50$, $3$, $10$ per sort, $61.5$ for the whole,
and the effective number of sorts is $1.37$. Three sorts are present but
the atom reads as about $1.4$, because color carries $85\%$ of the
strength.
\end{itemize}
\emph{Status: R41--R43; the v19 per-sort line is refined above.}

\section{Merged directly, without an ideas entry}

These went straight into the paper at your direction and are already
in \texttt{Pattern\_\allowbreak Arithmetic\_\allowbreak v15\_\allowbreak 22.tex}: the Pāṇini expansion
(pratyāhāra and vipratiṣedha as structural parallels, with the
Ingerman 1967 attribution of ``Pāṇini--Backus form'' and the
oversimplification caveat); the paragraph on summing separate
realities; the quantum-mechanics/general-relativity paragraph,
deliberately hedged so that it does not claim Pattern Arithmetic
bridges them; and the $N$/$D$ addition (filed afterwards as R5).

\section{What remains, in dependency order}

\begin{enumerate}
\item \textbf{Import the meta-vector taxonomy into Paper~1 (R15).}
It unblocks R14, which closes Gap~2 and gives location, the
timestamp and expiry registers, and $V_{\text{derived}}$ a home.
\item \textbf{Then merge the v17 resolution (R14)}, replacing the
paper's ``Open: location as a sort'' paragraph. Decide R19 (expired
contributions) at the same time or defer it explicitly. If R38 is
adopted, v17's ``per vector'' registers should read ``per core.''
\item \textbf{Decide the scope questions (R27, R28)}: Axiom~7's
metrics and the costumes, and whether to split. These also decide
whether R12, R25, and R41--R43 still matter for Paper~1.
\item \textbf{The atom's structure (R20--R22, R36--R40).} The word
is settled (R38: $N$ cores, $Q$ vectors). Decide the general atom form --- cores,
per-vector $\tau_i$, contribution size free (R36, R37) --- which
closes R20 in the same step. R21 and R22 can each be closed by a
sentence, whether a stated refusal or a stated rule, and the commuting
test (R39, R40) can then be run.
\item \textbf{Package the Inference-bound items (R17, R30--R35)}
as one coherent addition to the Inference paper: the candlestick,
the technique-selection argument, the spherical Douglas--Peucker
adaptation, the sphere-packing distance, and meta-vector grouping.
\item \textbf{Small residuals}: the one-sentence R26 clause, the
review's presentational items (R29), shell-indexed energy (R23), the
antipodal special case (R24), and whether $S_{\text{atom}}/N$ earns a
name.
\end{enumerate}

\end{document}
